\chapter{Objetivos}

\section{Objetivo Geral}

Estudar a detecção de vulnerabilidades em código-fonte com modelos de linguagem, considerando ao mesmo tempo a qualidade da detecção sob avaliação descontaminada e o custo computacional de técnicas de eficiência, como a quantização e o ajuste fino com QLoRA.

\section{Objetivos Específicos}

Os objetivos específicos do estudo compreenderam:

\begin{alineas}
    \item compreender a formulação da detecção de vulnerabilidades por função e as formas de uso de modelos de linguagem de código nessa tarefa;
    \item estudar os problemas de qualidade dos conjuntos de dados e o protocolo de avaliação do PrimeVul;
    \item estudar e implementar as métricas de detecção: F1, precisão, revocação, taxa de falsos positivos, VD-S e avaliação pareada;
    \item implementar o ajuste fino com QLoRA e cabeça de classificação nos modelos Qwen2.5-Coder de 3B e 7B;
    \item estudar os métodos de quantização NF4, AWQ e GPTQ e avaliar experimentalmente a implantação em NF4;
    \item medir o custo computacional das configurações avaliadas;
    \item registrar os resultados em relatório técnico e em artigo científico.
\end{alineas}
